% !TeX root = ../../main.tex
\section{分层筛选的执行示例与实现边界}
\subsection{候选重算的校正能力与覆盖边界}
为了直观区分分数恢复与候选恢复，表\ref{tab:acc-toy}给出一个人为构造的八行示例。目标为 Top-2，预筛保留3 行；表中分数仅用于解释机制，不来自论文实验或真实量化器测量。真实参考集合为行1 和行2，但行2 的近似分数落到第4 位，未进入预筛集合。

\begin{table}[htbp]\centering\caption{分层筛选的分析性示例：预筛遗漏不能由重算恢复}\label{tab:acc-toy}
\begin{tabular}{cccc}\toprule
行号&FP32 参考分数&近似分数&是否进入 Top-3 预筛\\\midrule
1&0.90&0.89&是\\
2&0.89&0.83&否\\
3&0.88&0.88&是\\
4&0.87&0.87&是\\
5&0.40&0.41&否\\
6&0.30&0.29&否\\
7&0.20&0.21&否\\
8&0.10&0.10&否\\\bottomrule
\end{tabular}\end{table}

对预筛集合 $\{1,3,4\}$执行高精度重算，得到最终结果 $\{1,3\}$，Recall@2 为 $1/2$。结果中每行的分数都可以是正确的 FP32 分数，但结果集合仍不完整。如果把预筛规模提高到4，则该示例可以恢复真实 Top-2。该现象说明，候选容量是质量补偿机制，其有效性取决于分数间隔和扰动大小。

另一个重要情况是近似阶段把行3 排在行1 之前，而二者都在候选中。此时 FP32 重算可以纠正相对顺序。因此，分层执行有意允许候选内部的排序误差，将高精度预算用于最终集合中的次序与分数。它要求保护的是候选覆盖，而不是让低比特阶段复现全部高精度排序。这个目标比逐项精确计算更宽松，也是节省计算的依据。

\subsection{局部候选容量的独立影响}
原架构中每个 Q-core 维护有限局部候选，再由后续模块处理汇集结果。局部选择与全局选择并不等价。假设两个分区各保留1 行，全局目标为 Top-2，分区 A 的分数为0.9 和0.8，分区 B 为0.3 和0.2。即便所有分数均精确，汇集后也只有0.9 和0.3，真实第二名0.8 已被删除。

该反例说明，局部选择的质量依赖分区容量及重要候选的分布。原方案使用32 个 Q-core 和每核16 项，产生最多512 个候选；“最多512 项”描述容量总和，不意味着获得了全局分数最高的512 行。实验中较高覆盖率来自具体数据、划分与容量配置，不能由 $32\times16$ 这一算式推导出保证。

可用每个分区真实 Top-K 数量的分布评估容量压力。设第 $p$ 个分区包含 $k_p$ 个真实目标，局部容量为 $l_p$，若所有分数精确，且次级排序规则一致，则 $l_p\geq k_p$ 对保存该分区目标是充分条件。实际低比特分数还会引入其他行对容量的竞争，所以条件需要进一步考虑近似排序，而非只统计精确目标数量。

为区分两种风险，离线分析可以保存分区编号、真实排名、近似排名和是否入选四类信息。某行若在精确局部 Top-$l_p$中却未进入近似局部集合，主要受数值扰动影响；若在精确局部选择中也被删除，则容量或分区是直接原因。该分解能够指导扩大候选、改变映射或调整剪枝，而不是统一提高所有精度。

\subsection{非零剪枝与行长度变化}
ReSparse 按元素幅值保留部分非零项，其直接目标是减少扫描记录。然而，不同行的原始非零数不同，相同保留比例不一定对应相同绝对长度。例如一行原有4 项，另一行原有40 项，保留一半后分别剩2 项和20 项；长行仍可能跨越多个传输包，短行则受边界信息和对齐成本影响更大。

剪枝后的行还可能变为空行。格式必须保留该行在原矩阵中的身份，否则候选行号与高精度副本不能对应。对空行，数学得分为零，但如果存在负分数，零分行可能影响 Top-K 结果。因此，不能无条件把空行从候选语义中删除。实现是否需要显式处理它们取决于输入约束和输出规则，其适用性由数据条件与输出规则共同决定。

对于含正负值的向量，删除小幅值元素并不保证得分单调下降。若被删项原本贡献负值，近似得分反而增大；若正负项存在抵消，较小的单项变化也可能导致较大的相对误差。因此，幅值剪枝的合理性应通过目标候选覆盖验证，不能把“删除小值”直接解释为“只删除无贡献元素”。

量化顺序也会影响结果。先剪枝再按新最大幅值缩放，与先量化再删除小项可能产生不同的整数分布。第3 章实现应依据原方案的预处理次序复现，同时保留原始矩阵作为 FP32 重算与质量参考。参数校准必须固定这个次序，否则同一保留比例和同一位宽未必对应同一候选集合。

\subsection{流式执行中的依赖与反压}
Q-core 的高并行度只有在数据持续到达、结果能够持续消费时才能转化为吞吐。稀疏行累加依赖行边界信息，Top-K 维护又依赖完整行分数。若排序模块处理速度不足，其输入队列会填满，反压会沿流水返回乘加和读取模块。此时单独增加乘法器数量不能提高完整查询速率。

类似地，多个 Q-core 汇集候选时，汇合点可能形成瞬时突发。共享 F-core 需要根据候选行号读取高精度行，如果某次查询候选行特别长，精排阶段服务时间会增加。合理的流水分析既要看平均候选数量，也要看候选非零总数与请求长度分布。512 个短行与512 个长行可能具有明显不同的重算成本。

可将每条查询的候选重算工作量写为
\begin{equation}
 z_{\mathrm{re}}(q)=\sum_{i\in\mathcal C(q)}\nnz(A_{i,:}),
\end{equation}
其中 $\mathcal C(q)$ 为实际入选候选集合。若使用剪枝后的行长度估计此量，会低估读取原始 FP32 副本的代价；两阶段使用不同矩阵表示，计量时必须明确引用哪份数据。图中报告的总体节省应与这个实际访问集合相联系。

\subsection{吞吐与单查询延迟的区别}
对只有一个请求的场景，查询要顺序经过预筛与精排，完成时间接近两阶段时间之和。对连续请求，若缓冲和数据依赖允许重叠，预筛第 $q+1$ 条查询可以与精排第 $q$ 条查询并行，此时稳态吞吐由较慢阶段决定。共享内存端口、查询缓存复用和结果缓冲容量都可能限制这种重叠。

因此，实验的 QPS 结果需要结合请求提交方式理解。若基线采用批量 GPU 处理，而 FPGA 按流式查询统计，应该确认两者都包含同样的结果生成步骤，并说明批量对延迟的影响。高 QPS 并不自动意味着相同的低单查询延迟，也不意味着所有到达率下的服务尾延迟都会改善。

上述分析刻画了流水执行与查询到达方式的关系，服务级性能仍需在对应负载下验证。后续测量可记录低负载延迟、饱和吞吐和不同到达率下的排队时间，分别判断架构适合交互查询还是批量检索。这样能够把计算能力与服务能力区分开，而不改变本章已完成的硬件验证结论。

\subsection{参数选择与实现可迁移性}
位宽、剪枝比例和候选容量形成联合配置空间。较低位宽可能需要更大候选集合；更强剪枝也可能增加精排量，因而不应分别选择各参数的最优值后直接组合。较稳妥的方法是在固定验证查询上筛除不满足质量目标的配置，再比较这些配置的完整查询时间和资源占用。

在迁移到另一种 FPGA 或不同 HBM 接口时，INT6 配置未必仍最优。数据包位宽、DSP 支持的乘法形状、片上存储粒度和目标频率都会改变资源收益。如果某平台原生支持 INT8 向量指令而不支持 INT6，位级压缩还可能需要解包代价。本文结果支撑的是特定可定制硬件条件下的选择，而非所有处理器上的位宽优越性。

维度增加也影响查询缓存和索引编码。若列索引超过原字段范围，需要扩大索引或引入分块，前述打包密度随之变化。若查询无法完整复制到每个计算单元，访存争用也可能成为新瓶颈。这说明分层思想具有可迁移性，但具体字段、并行度和流水参数仍需要重新设计与验证。
